\documentclass[11pt,a4paper]{article}

\usepackage[T1]{fontenc}
\usepackage[utf8]{inputenc}
\usepackage{lmodern}
\usepackage[a4paper,margin=1in]{geometry}
\usepackage{booktabs}
\usepackage{longtable}
\usepackage{tabularx}
\usepackage{array}
\usepackage{enumitem}
\usepackage{hyperref}
\usepackage{microtype}
\usepackage{amsmath}
\usepackage{amssymb}
\usepackage{xcolor}

\hypersetup{
  colorlinks=true,
  linkcolor=blue!50!black,
  urlcolor=blue!50!black,
  citecolor=blue!50!black,
  pdftitle={NewsProp Research Paper},
  pdfauthor={NewsProp Five-Person Team}
}

\setlist[itemize]{noitemsep,topsep=2pt,leftmargin=1.2em}
\setlist[enumerate]{noitemsep,topsep=2pt,leftmargin=1.5em}

\title{\textbf{Modeling the Propagation of Fake News in a Synthetic Moldovan Information Network}}
\author{
  Team Member 1 \and Team Member 2 \and Team Member 3 \and Team Member 4 \and Team Member 5
}
\date{\today}

\begin{document}
\maketitle

\begin{abstract}
NewsProp is a reproducible simulation framework for studying misinformation spread in a Moldova-inspired information environment. The system combines scraped fake-news, mainstream-news, and Telegram-style social content with a multilingual normalization pipeline, a synthetic scale-free social graph, and a SEIZ-style agent-based diffusion model. The current verified dataset pipeline processed 1,121 usable records after quality filtering. On a 1,000-node calibrated social network with 2,991 edges and 50 identified hub nodes, the baseline simulation produced a mean peak infected population of 886.3 agents with a 95\% bootstrap confidence interval of [843.20, 929.01]. When a pre-bunking intervention was applied to hub nodes at the start of the simulation, peak infection dropped to 818.1 agents with a 95\% bootstrap confidence interval of [784.39, 852.20], a reduction of 7.69\%, while infected area under the curve decreased by 8.68\%. A sensitivity run using an expanded Telegram scrape kept 2,700 normalized records and produced a smaller but still positive intervention effect. These results suggest that targeted inoculation of influential nodes can reduce misinformation diffusion in synthetic Moldova-like networks, while corpus choice affects effect magnitude.
\end{abstract}

\section{Introduction}
Online misinformation behaves like a contagious process: it spreads through social ties, adapts through resharing, and affects people differently depending on their trust, skepticism, and exposure patterns. Moldova is a particularly relevant context for studying this problem because its media ecosystem contains a mix of mainstream news consumption, social-media dependence, and politically charged narratives that move across multiple platforms.

The goal of NewsProp is to create a scientifically interpretable simulation rather than a generic demo. The project therefore joins data collection, multilingual preprocessing, synthetic population generation, and intervention testing into one reproducible workflow. The central research question is:

\begin{quote}
How does fake news propagate through a synthetic Moldovan information network, and to what extent can a targeted pre-bunking strategy reduce its spread?
\end{quote}

\section{System Overview}
NewsProp is organized into five phases:

\begin{enumerate}
  \item Data collection and scraping.
  \item NLP normalization and impact scoring.
  \item Synthetic network generation.
  \item SEIZ-style propagation simulation.
  \item Analysis, visualization, and paper generation.
\end{enumerate}

This phase structure matters because the later claims depend on artifacts generated by earlier phases, not on hand-crafted examples.

\section{Data Pipeline}
The current repository includes scrapers and datasets for:

\begin{itemize}
  \item fake news and debunked claims from StopFals and related sources,
  \item real news from Moldovan mainstream sources such as \texttt{stiri.md},
  \item social-channel content from Telegram-focused scraping work.
\end{itemize}

The Phase 2 normalization pipeline standardizes these sources into a common schema and computes downstream features such as:

\begin{itemize}
  \item detected language,
  \item English translation when enabled,
  \item VADER sentiment,
  \item emotional score,
  \item emotional intensity,
  \item engagement totals,
  \item optional sentence-transformer embeddings.
\end{itemize}

For the currently verified execution path used in this draft, Phase 2 was run in a robust smoke configuration with translation and embeddings disabled:

\begin{verbatim}
.\.venv\Scripts\python.exe phase2\pipeline.py `
  --disable-translation `
  --disable-embeddings `
  --output-dir phase2\outputs\smoke
\end{verbatim}

This run kept 1,121 records and rejected 1 record during quality filtering.

\section{Network Construction}
Phase 3 generates a Moldova-inspired synthetic network using a Barabasi-Albert scale-free graph. Each node is assigned agent attributes such as age, critical thinking, media trust, and preferred channel.

Hub nodes are identified through centrality-based ranking. In the current verified 1,000-node run:

\begin{itemize}
  \item nodes: 1,000,
  \item edges: 2,991,
  \item hub nodes: 50,
  \item hub degree share: 24.86\%,
  \item average degree: 5.982,
  \item degree Gini: 0.3796.
\end{itemize}

These values support the intuition that a relatively small number of influential nodes account for a disproportionate share of network connectivity.

\section{Simulation Method}
Phase 4 uses a SEIZ-like model with the following states:

\begin{itemize}
  \item \texttt{S}: Susceptible,
  \item \texttt{E}: Exposed,
  \item \texttt{I}: Infected,
  \item \texttt{Z}: Skeptic.
\end{itemize}

Agents interact across their graph edges and through preferred channels such as social diffusion and broadcast-like exposure. The simulator also includes a small distortion effect that perturbs news parameters during resharing. This approximates the way messages mutate as they spread.

The simulator now prefers Phase 2-normalized outputs when available. That means Phase 4/5 can directly consume impact score, Phase 2 emotional intensity, and cleaned engagement metrics instead of reconstructing all news features from raw scraped data.

\subsection{Reproducible Execution Path}
The current reproducible workflow is intentionally simple:

\begin{enumerate}
  \item bootstrap a local virtual environment and install dependencies,
  \item run Phase 2 normalization,
  \item run Phase 3 network generation,
  \item run Phase 4/5 simulation on the Phase 3 network while consuming the Phase 2 normalized dataset.
\end{enumerate}

This workflow is encoded in \texttt{scripts/run\_full\_pipeline.ps1}. For smaller verification runs, the same logic can be executed with reduced node counts, run counts, and tick counts.

\subsection{Current Parameterization}
The verified full experiment used the following settings:

\begin{table}[h]
\centering
\begin{tabular}{lr}
\toprule
Parameter & Value \\
\midrule
agents & 1000 \\
BA attachment parameter \textit{m} & 3 \\
hub fraction & 0.05 \\
simulation runs per scenario & 10 \\
ticks per run & 80 \\
initial infected seeds & 8 \\
random seed & 42 \\
\bottomrule
\end{tabular}
\caption{Verified baseline simulation settings.}
\end{table}

\section{Experimental Design}
The primary verified experiment compares two scenarios:

\begin{enumerate}
  \item Baseline spread.
  \item Pre-bunking intervention where hub nodes are initialized as skeptics.
\end{enumerate}

The verified full run used 1,000 agents, 10 runs per scenario, 80 ticks, 8 initial infected seeds, a 5\% hub fraction, and random seed 42.

The news item selected by the simulator in the verified full run was a fake-news article from Facebook, with Phase 2-derived properties:

\begin{itemize}
  \item controversy score: 0.8225,
  \item emotional intensity: 0.9965,
  \item impact score: 0.9965,
  \item transmission probability: 0.7748.
\end{itemize}

\section{Current Results}
The verified full experiment produced the following summary:

\begin{table}[h]
\centering
\begin{tabularx}{\linewidth}{lrrX}
\toprule
Metric & Baseline & Pre-bunked hubs & Relative change \\
\midrule
Peak infected & 886.3 & 818.1 & -7.69\% \\
Peak infected 95\% bootstrap CI & [843.20, 929.01] & [784.39, 852.20] & lower under intervention \\
Tick of peak & 6 & 7 & delayed by 1 tick \\
Infected AUC & 67,971.95 & 62,069.85 & -8.68\% \\
Infected AUC 95\% bootstrap CI & [64,735.38, 70,838.36] & [59,425.31, 64,719.04] & lower under intervention \\
Final infected & 886.3 & 818.1 & reduced \\
\bottomrule
\end{tabularx}
\caption{Scenario-level results for the verified full experiment.}
\end{table}

The intervention reduced peak infected by a paired mean of 68.2 agents with a 95\% bootstrap confidence interval of [25.49, 116.61], reduced infected AUC by a paired mean of 5902.1 with a 95\% bootstrap confidence interval of [2612.17, 9420.55], and delayed the peak by a mean of 0.9 ticks with a 95\% bootstrap confidence interval of [0.6, 1.2]. The intervention improved peak infected and infected AUC in 80\% of paired runs, and it delayed or preserved the peak timing in 100\% of paired runs.

\subsection{Telegram Refresh Sensitivity}
The larger Telegram scrape from the Lucian-authored branch changed the experiment in a meaningful way. After normalizing 2,700 records and rerunning the simulation on the same 1,000-node network, the intervention still reduced average spread, but the effect became smaller and less certain for some metrics:

\begin{itemize}
  \item paired mean peak infected reduction: 24.6 agents,
  \item paired mean infected AUC reduction: 3,272.6,
  \item paired mean peak delay: 3.1 ticks,
  \item positive intervention share: 70\% for peak infected and infected AUC,
  \item positive delay share: 100\% for peak timing.
\end{itemize}

The bootstrap confidence intervals for peak infected and infected AUC cross zero in this sensitivity run, which is scientifically useful because it shows the effect is not uniformly strong under all data choices. Instead, the intervention appears robust in direction, but its magnitude depends on the news corpus used to define the dominant misinformation item.

\section{Interpretation}
The most important current finding is not just that pre-bunking helps, but that it helps in a network-aware way. The intervention does not need to immunize the whole population. Instead, it reduces spread by focusing on the structurally influential part of the network. That aligns with the project's intended thesis: misinformation control is more efficient when it is targeted rather than uniform.

The bootstrap intervals strengthen this claim by showing that the reduction is not only visible in point estimates, but also persists when run-to-run variability is summarized statistically. The paired-run results add a second layer of evidence: the intervention is not merely better on average, but is usually better for matched runs and always delays or preserves the timing of the peak in the current verified experiment.

The Telegram refresh sensitivity result adds an important caveat: broader source coverage does not necessarily amplify the effect size. Instead, it can reveal that the apparent strength of the intervention depends on which specific news item the corpus elevates as the most transmissible example.

The sampled translation comparison adds a second caveat. On a deterministic 200-record subset of the refreshed Telegram corpus, translation changed the Phase 2 sentiment calibration sharply: the translation-enabled run produced a global mean sentiment of -0.0782, while the translation-disabled run produced 0.6278. When both Phase 2 outputs were fed into the same 200-node smoke network, the translation-enabled path selected a different dominant fake-news article and slightly increased the intervention effect, with peak reduction rising from 1.65\% to 2.69\% and infected AUC reduction rising from 2.98\% to 4.07\%. This makes translation a scientifically important preprocessing choice, but it also means the paper should describe the setting explicitly rather than treat it as a hidden implementation detail.

\section{Limitations}
The present draft still has important limitations:

\begin{itemize}
  \item the main verified run used a translation-disabled Phase 2 configuration for speed and robustness,
  \item uncertainty intervals for intervention outcomes still need to be expanded into a broader experiment matrix,
  \item the network is synthetic rather than directly estimated from Moldova-specific social graph data,
  \item the current simulation parameters remain partially heuristic,
  \item literature support and formal related-work citations are still pending,
  \item the paired-delta analysis code has been verified, but the manuscript still needs a Tectonic build path on this machine.
\end{itemize}

\section{Next Steps}
The next scientific improvements should focus on:

\begin{enumerate}
  \item comparing translation-disabled and translation-enabled Phase 2 feature paths,
  \item refreshing scraper documentation and dataset provenance,
  \item adding Moldova-specific literature and official statistics to support the assumptions,
  \item refining the final academic formatting and bibliography,
  \item expanding from a single verified experiment into a broader experiment matrix.
\end{enumerate}

\section*{UTM Guidance}
The manuscript follows the UTM writing guidance and scientific-article expectations summarized from the official repository and journal resources:

\begin{itemize}
  \item the thesis-writing guide emphasizes structured scientific writing, methodological correctness, and coherent academic style \cite{utm_thesis_guide},
  \item the UTM journal instructions emphasize author guidelines, manuscript preparation discipline, and clear submission requirements \cite{utm_jes_instructions}.
\end{itemize}

\section*{Appendix A: Planned Artifact Table}
The final paper should reference the following generated artifacts directly:

\begin{table}[h]
\centering
\begin{tabularx}{\linewidth}{lX}
\toprule
Artifact & Role in paper \\
\midrule
\texttt{outputs/phase3\_network/network\_summary.json} & network and calibration summary \\
\texttt{outputs/phase45/results\_summary.json} & top-level experiment results \\
\texttt{outputs/phase45/run\_metrics.csv} & per-run scenario outcomes \\
\texttt{outputs/phase45/paired\_run\_comparison.csv} & matched baseline vs intervention comparison \\
\texttt{outputs/phase45/infected\_curve\_comparison.png} & primary intervention figure \\
\texttt{outputs/phase45/state\_trajectories.png} & state evolution visualization \\
\texttt{outputs/phase45/network\_snapshot.png} & qualitative network illustration \\
\bottomrule
\end{tabularx}
\caption{Artifacts to cite in the final submission.}
\end{table}

\section*{References}
\begin{thebibliography}{99}
\sloppy
\bibitem{utm_thesis_guide}
J. Ciumac and L. Gherciu-Mustea\c{t}\u{a}, \emph{Ghid de redactare \c{s}i sus\c{t}inere a tezei de licen\c{t}\u{a}}, Universitatea Tehnic\u{a} a Moldovei, 2010. Available at \href{https://repository.utm.md/handle/5014/15847}{UTM repository entry}.

\bibitem{utm_jes_instructions}
\emph{Instructions for authors -- Journal of Engineering Science}, Technical University of Moldova. Available at \href{https://jes.utm.md/instructions-for-authors/}{JES author instructions}.

\bibitem{networkx_ba}
NetworkX Documentation, \emph{barabasi\_albert\_graph}. Available at \href{https://networkx.org/documentation/stable/reference/generated/networkx.generators.random_graphs.barabasi_albert_graph.html}{NetworkX BA graph docs}.

\bibitem{networkx_betweenness}
NetworkX Documentation, \emph{betweenness\_centrality}. Available at \href{https://networkx.org/documentation/stable/reference/algorithms/generated/networkx.algorithms.centrality.betweenness_centrality.html}{NetworkX betweenness docs}.

\bibitem{networkx_eigenvector}
NetworkX Documentation, \emph{eigenvector\_centrality}. Available at \href{https://networkx.org/documentation/stable/reference/algorithms/generated/networkx.algorithms.centrality.eigenvector_centrality.html}{NetworkX eigenvector docs}.

\bibitem{worldbank}
World Bank API indicators for Moldova: SP.POP.1564.TO.ZS, SP.POP.65UP.TO.ZS, IT.NET.USER.ZS, SP.URB.TOTL.IN.ZS. Available via \href{https://api.worldbank.org/}{World Bank API}.

\bibitem{ijc}
IJC media audience study article on Moldova media consumption. Available at \href{https://cji.md/en/media-audience-study-launched-by-the-ijc-social-networks-have-become-the-main-source-of-information-for-media-consumers-in-the-republic-of-moldova/}{IJC study article}.
\end{thebibliography}

\end{document}
